% ============================================================
%  第5章 OpenCL运行时和并发模型
%  源文件：content/chapter5/
% ============================================================

支持OpenCL设备已经相当多了，从离散式显卡(其拥有成千上万个“核芯”)，到嵌入式设备上的中央处理器(CPU)。能获得如此广泛的支持，很重要的原因在于OpenCL的内存模型和执行模型。OpenCL通过抽象模型的方式跨越硬件架构的鸿沟(不需要遵循特定硬件上的编程规则)，同时让编程者在相应的硬件上获得更优的性能，且不用代码做太多修改。本章，我们会讨论执行模型中的不同组件。


% ---- 5.1-chinese.md ----
\section{命令和排队模型}

OpenCL基于任务并行/主机控制模型，其每个任务均为数据并行。其通过使用每个设备上线程安全的命令队列作为该模型的保证。内核、数据搬移，以及其他操作并非是使用者调用一些运行时函数，就能简单的执行。这些操作使用异步入队操作，将命令送入指定队列，并在未来的某个时刻执行该命令。执行的同步点在于主机端的命令队列和设备端的命令队列。

OpenCL命令队列中命令可以是，执行内核、内存数据转移或是同步命令。要想了解命令的执行结果，只能等到命令队列到达同步点才能看到。下面列举几个主要的同步点：

\begin{itemize}
\item 等待指定 OpenCL 事件的完成
\item 调用\texttt{clFinish()}函数，将阻塞主机端的线程，直到命令队列中的命令全部执行完毕
\item 执行一次阻塞式内存操作
\end{itemize}

\subsection{阻塞式内存操作}

阻塞式内存操作应该是最常见，而且是最简单的同步方式。这种操作会阻塞主线程的进行，直到内存数据传输完成。数据传输API都具有一个参数，可以决定是否使用阻塞的方式进行数据传输。\texttt{clEnqueueReadBuffer()}中的blocking\_\allowbreak{}read就是用来控制是否使用阻塞方式进行数据传输。

\begin{codebox}[cpp]
cl_int
clEnqueueReadBuffer(
  cl_command_queue command_queue,
  cl_mem buffer,
  cl_bool blocking_read,
  size_t offset,
  size_t size,
  const void *ptr,
  cl_uint num_events_in_wait_list,
  const cl_event *event_wait_list,
  cl_event *event)
\end{codebox}

将数据从设备端获取，或是传输到设备端，通常都会使用同步的方式进行内存操作。例如，从设备端获取数据到主机端时，在传输没有完成前，主机是不能访问其数据，否则将会导致一些未定义行为。因此，blocking\_\allowbreak{}read参数可以设置成CL\_\allowbreak{}TRUE，用以阻塞主线程，直到数据传输完成主线程才能继续进行下面的操作。使用这种方式进行同步，可以直接获取数据，之后就不需要在进行额外的同步了。阻塞和非阻塞式内存操作将在第6章进行详细讨论。

\subsection{事件}

这里先回顾一下第3章中，事件可以用来指定命令之间的依赖关系。每个\texttt{clEnqueue*}API都可以产生一个与其入队命令相关的事件对象，我们可以通过事件对象来查询命令的状态，并且指定对应命令所要依赖命令所对应的事件队列。生成的事件可以作为一种依赖机制，以便OpenCL运行时实现其所要执行的任务图[译者注1]。

随着命令入队和出队，以及命令的执行，事件都会持续的更新其状态。命令状态列举如下：

\begin{itemize}
\item 已入队(Queued)：该命令已经在命令队列中占据了一席之地
\item 已提交(Submitted)：该命令已经从命令队列中移除，并提交给设备执行
\item 已就绪(Ready)：该命令已经准备好在设备端执行
\item 已运行(Running)：该命令已经在设备端执行，但并未完成
\item 已完成(Ended)：该命令已经在设备端执行完成
\item 已结束(Complete)：该命令以及其子命令都已经执行完成
\end{itemize}

因为异步是OpenCL天生的特性，所以其API不能简单的返回错误码，或与命令执行相关的性能数据(可以通过API提供的参数获取入队时的错误码，以及入队的时间点)。不过，OpenCL也提供查询命令相关错误码的机制。要查询对应的错误码，需要借助与该命令关联的事件对象。事实上，完成本身也可以视为与其他状态类似的一种状态。查询事件对象的状态，需要调用\texttt{clGetEventInfo()}API，将CL\_\allowbreak{}EVENT\_\allowbreak{}COMMAND\_\allowbreak{}STATUS\_\allowbreak{}EXECUTION\_\allowbreak{}STATUS传递给param\_\allowbreak{}name参数即可。

\begin{codebox}[cpp]
cl_int
clGetEventInfo(
  cl_event event,
  cl_event_info param_name,
  size_t param_value_size，
  void *param_value,
  size_t *param_value_size_ret)
\end{codebox}

当命令完全正确无误完成，相关事件对象的状态会设置成CL\_\allowbreak{}COMPLETE。注意“完全”值的是其命令本身，以及其相关子命令，都需要正确完成。子命令指的是内核入队其子内核，子内核的运行称为子命令。有关设备端入队的内容会在之后(本章)进行讨论。

当命令非正常结束，为正常完成时，相关事件对象的状态通常会返回一个负值。这种情况下，当有命令异常终止，那么使用同样上下文对象的命令队列则不再可用，相关命令队列上的命令将会全部退出。

调用\texttt{clWaitForEvents()}将会阻塞主机端的执行，直到指定的事件对象链表上相关的命令全部解除，才会解除对主机主线程的阻塞。

\begin{codebox}[cpp]
cl_int
clWaitForEvent(
  cl_uint num_events,
  const cl_event *event_list)
\end{codebox}

\subsection{命令栅栏和命令标识}

想要不阻塞主机端线程，用于同步命令的方式也存在——入队一个命令栅栏。命令栅栏从原理来说，与在主机端调用\texttt{clWaitForEvent()}类似，但这种方式不会显式的进行处理，而是运行时库会内部处理。入队栅栏需要使用\texttt{clEnqueueBarrierWithWaitList()}API[译者注2]，可将事件列表作为可选参数传入。如果没有提供事件，栅栏会等待之前入队的所有命令完成(后续入队的命令将不会执行，直到前面所有的命令执行完成)。

标识(Marker)的作用于栅栏类似，其使用\texttt{clEnqueueMarkerWithWaitList()}API[译者注2]进行入队。栅栏和标识的区别在于，标识不会阻塞之后入队的命令。因此，当设备完成所有指定事件时，标识就允许编程者去查询指定事件状态，而不会阻碍其中一些命令的执行。

使用这两种同步命令，以及事件的方式，就给OpenCL提供了更多的可能，能够完成更加复杂的任务图，完成高难度的动作。这种方式在使用乱序命令队列时尤为重要，其能在运行时优化命令的调度。

\subsection{事件回调}

OpenCL允许用户为事件对象注册回调函数。当事件对象到达某一指定状态时，回调函数便会调用。\texttt{clSetEventCallback()}函数就可为OpenCL的事件对象注册回调函数：

\begin{codebox}[cpp]
cl_int
clSetEventCallback(
  cl_event event,
  cl_int command_exec_callback_type,
  void (CL_CALLBACK *pfn_event_notify)(
    cl_event event,
    cl_int event_command_exec_status,
    void *user_data),
  void *user_data)
\end{codebox}

command\_\allowbreak{}exec\_\allowbreak{}callback\_\allowbreak{}type用来指定回调函数在何时调用。可能的参数只有：CL\_\allowbreak{}SUBMITTED，CL\_\allowbreak{}RUNNING和CL\_\allowbreak{}COMPLETE。

关于回调函数调用顺序的保证是有限的。举个例子，不同的回调函数注册了事件状态为CL\_\allowbreak{}SUBMITTED和CL\_\allowbreak{}RUNNING，当事件对象的状态发生变化就能确保按照正确的顺序进行，而回调函数的调用顺序就无法保证正确。我们将在第6章和第7章讨论到，内存状态只能在事件对象为Cl\_\allowbreak{}COMPLETE是确定(事件对象为其他状态时，无法确定内存状态)。

\subsection{使用事件分析性能}

确定一个命令的执行时间时，会将对应的数值传递给事件对象，命令在不同的状态时都由对应的计时器进行计时。不过，这样的计时需要开启命令队列的计时功能，需要在创建命令队列时，将CL\_\allowbreak{}QUEUE\_\allowbreak{}PROFILING\_\allowbreak{}ENABLE加入properties参数内，提供给\texttt{clCreateCommandQueueWithProperties()}。

任何与命令相关的时间对象，都可以通过调用\texttt{clGetEventProfilingInof()}获取性能信息：

\begin{codebox}[cpp]
cl_int
clGetEventProfilingInfo(
  cl_event event,
  cl_profiling_info param_name,
  size_t param_value_size,
  void *param_value,
  size_t *param_value_size_ret)
\end{codebox}

通过对不同状态的耗时查询，编程者可以知道命令每个状态在队列中的耗时情况。同样，编程者通常就想知道对应命令实际运行的时间(比如：数据传输的耗时情况，内核执行的耗时情况)。为了确定命令整体耗时情况，可以将CL\_\allowbreak{}PROFILING\_\allowbreak{}COMMAND\_\allowbreak{}START和CL\_\allowbreak{}PROFILING\_\allowbreak{}COMMAND\_\allowbreak{}END作为实参传入param\_\allowbreak{}name中。如果需要将子内核执行的时间计算在内，就需要传递CL\_\allowbreak{}PROFILING\_\allowbreak{}COMMAND\_\allowbreak{}COMPLETE作为实参。OpenCL定义的计时器，其计时的精度必须是纳秒级别。

\subsection{用户事件}

目前我们所见到的事件对象都是入队命令API，通过其参数中提供一个事件对象指针的方式所产生的事件对象。不过，当编程者想要OpenCL命令等待一个主机端的事件，应该怎么办呢？例如，编程者可能需要OpenCL进行数据传输，这个数据传输任务需要等待某个文件更新之后再进行传输。为了应对这种情况，OpenCL 1.2标准中添加了用户事件(user event)。

\begin{codebox}[cpp]
cl_event
clCreateUserEvent(
  cl_context context,
  cl_int *errcode_ret)
\end{codebox}

这样创建的事件对象，其状态就是由应用开发者来决定，而非OpenCL运行时来决定。用户事件的状态数量会有限制，用户事件可以为“已提交”状态(CL\_\allowbreak{}SUBMITTED)，也可为“已完成”(CL\_\allowbreak{}COMPLETE)，或是某一个错误状态。当一个用户事件对象创建，其执行状态就被设置成CL\_\allowbreak{}SUBMITTED。

用户事件的状态可以通过\texttt{clSetUeserEventStatus()}进行设置：

\begin{codebox}[cpp]
cl_int clSetUserEventStatus(
  cl_event event,
  cl_int execution_status)
\end{codebox}

execution\_\allowbreak{}status参数指定需要设置的新执行状态。用户事件对象可以将状态设置为两种，第一种可为CL\_\allowbreak{}COMPLETE；第二种可为一个负数，表示一个错误。设置成负值意味着所有等待该用户事件的已入队命令都会被终止。还有一点需要注意的是，\texttt{clSetUserEventStatus()}只能对一个用户事件对象使用一次，也就是对一个对象只能设置一次执行状态。

\subsection{乱序任务队列}

第3章和第4章中的OpenCL例子都是以顺序命令队列(默认方式)进行。顺序命令队列能保证命令将按入队的顺序，在对应的设备上执行。不过，OpenCL队列也可以乱序执行。一个乱序队列时没有什么执行顺序可言的。硬件支持直接内存访问(direct memory access, DMA)引擎可以与计算单元并行执行，或设备可以同时执行多个内核。像这样的硬件设备就完全可以支持OpenCL运行时使用乱序命令队列，其可以自由且并行的调度这些操作(就是无法保证哪个命令先执行)。

创建命令队列的API(\texttt{clCreateCommandQueueWithProperties()})可以设置乱序的标志位，这种方式到现在我们还没有尝试过。只要属性参数中包含CL\_\allowbreak{}QUEUE\_\allowbreak{}OUT\_\allowbreak{}OF\_\allowbreak{}ORDER\_\allowbreak{}EXEC\_\allowbreak{}MODE\_\allowbreak{}ENABLE就可以产生支持乱序执行的命令队列。需要注意的，支持乱序的队列需要创建在支持乱序执行的设备上。

代码清单5.1中，展示了使用乱序命令队列程序主机代码的一部分，其包括写入输入数组、执行两个内核，并且将数据读取到主机端。这一系列命令的顺序由特定的事件对象指定，事件对象中记录了各个命令的依赖关系。任务图就由这些事件对象创建，这样即使是乱序执行也能保证命令执行的顺序正确。其中内存传输函数是非阻塞的，并且最后的同步，是显式的在主机端对读取事件对象使用\texttt{clWaitForEvents()}完成的。使用非阻塞形式的内存传输需要格外注意，因为当使用乱序队列时，可能会有潜在的传输区域或执行区域的重叠。

\begin{codebox}[cpp]
// -----------------------------
// Relevant host program
// -----------------------------

// Create the command-queue
cl_command_queue_properties properties = CL_QUEUE_OUT_OF_ORDER_EXEC_MODE_ENABLE;
cl_command_queue queue = clCreateCommandQueueWithProperties(context, devices[0], &properties, NULL);

// Declare the events
cl_event writeEvent, kernelEvent0, kernelEvent1, readEvent;

// Create the buffers
cl_mem input = clCreateBuffer(context, CL_MEM_READ_ONLY, 32 * sizeof(float), NULL, NULL);
cl_mem intermediate = clCreateBuffer(context, CL_MEM_READ_WRITE, 32 * sizeof(float), NULL, NULL);
cl_mem outpt = clCreateBuffer(context, CL_MEM_WRITE_ONLY, 32 * sizeof(float), NULL, NULL);

// Write the input data
clEnqueueWriteBuffer(queue, input, CL_FALSE, 0, 32 * sizeof(float), (void *)hostInput, 0, NULL, &writeEvent);

// Set up the execution unit dimensions used by both kernels
size_t localws[1] = {8};
size_t globalws[1] = {32};

// Enqueue the first kernel
clSetKernelArg(kernel, 0, sizeof(cl_mem), (void *)&input);
clSetKernelArg(kernel, 1, sizeof(cl_mem), (void *)&intermediate);
clSetKernelArg(kernel, 2, 8 * sizeof(float), NULL);
clEnqueueNDRangeKernel(queue, kernel, 1, NULL, globalws, localws, 1, &kernelEvent0, &kernelEvent1);

// Read output data
EnqueueNDRangeKernel(queue, output, CL_FALSE, 0, 32 * sizeof(float), (void *)&hostOutput, 1, &kernelEvent1, &readEvent);

// Block until the read has completed
clWaitForEvents(1, &readEvent);

clReleaseEvent(writeEvent);
clReleaseEvent(kernelEvent);
clReleaseEvent(readEvent);
\end{codebox}

代码清单5.1 同时入队两个内核，使用事件来控制依赖关系

乱序命令队列并不保证其执行顺序就是乱序的。鲁棒性较好的应该避免死锁的发生，当意识到有死锁发生的可能，那最好的方式就是串行执行入队的命令。我们再来看一下代码清单5.1的代码，即便是是顺序执行，程序也能得到正确的结果。不过，当“写入输入数据”和“执行第一个内核”这两个任务在源码中颠倒，那么开发者还是期望队列是以乱序的方式执行。但是，当队列顺序执行的时候，就会产生死锁，因为内核执行事件将会等待写入操作的完成。

\myRefRule

\begin{myRefs}

[译者注1] ”任务图“的概念类似于OpenVX中Graph的概念，有兴趣的读者可以查阅OpenVX的官方文档。

[译者注2]这两个API很相似，不过其描述不同。下面引用OpenCL官方的描述：

\begin{myQuote}
clEnqueueBarrierWithWaitList: Enqueues a barrier command which waits for either a list of events to complete, or if the list is empty it waits for all commands previously enqueued in command\_\allowbreak{}queue to complete before it completes. This command blocks command execution, that is, any following commands enqueued after it do not execute until it completes. This command returns an event which can be waited on, i.e. this event can be waited on to insure that all events either in the event\_\allowbreak{}wait\_\allowbreak{}list or all previously enqueued commands, queued before this command to command\_\allowbreak{}queue, have completed.
\end{myQuote}

\begin{myQuote}
clEnqueueMarkerWithWaitList： Enqueues a marker command which waits for either a list of events to complete, or if the list is empty it waits for all commands previously enqueued in command\_\allowbreak{}queue to complete before it completes. This command returns an event which can be waited on, i.e. this event can be waited on to insure that all events either in the event\_\allowbreak{}wait\_\allowbreak{}list or all previously enqueued commands, queued before this command to command\_\allowbreak{}queue, have completed.
\end{myQuote}

\end{myRefs}



% ---- 5.2-chinese.md ----
\section{多命令队列}

如果系统里面有多个设备(比如，CPU和GPU，或多个GPU)，每一个设备都需要拥有自己的命令队列。不过，OpenCL允许将多个命令队列，通过同一个上下文对象映射到同一个设备上。这对重叠执行独立或重叠命令，以及主机和设备通讯十分有用，并且这种方式是替代乱序队列的一种策略。了解同步的能力，以及主机端和设备端的内存模型(第6章和第7章将详细讨论)对于管理多个命令队列来说是非常必要。

图5.1展示了一个OpenCL上下文中具有两个设备的情况。为不同的设备创建不同的命令队列。代码清单5.2中展示了创建两个命令队列的相关代码。这里尤其要注意的是，使用OpenCL事件对象进行同步时，只能针对同一上下文对象中的命令。如果是不同的上下文中的设备，那么事件对象的同步功能将会失效，并且要在这种情况下共享数据，就需要在两个设备之间进行显式的拷贝。

\myGraphic{0.85}{images/chapter5/5-1.png}{图5.1 在同一上下文中为不同设备创建的多个命令队列。图中为两个设备各创建了一个命令队列。}

\begin{codebox}[cpp]
// Obtain devices of both CPU and GPU types
cl_device_id devices[2];
err_code = clGetDeviceIDs(NULL, CL_DEVICE_TYPE_CPU, 1, &devices[0], NULL);
err_code = clGetDeviceIDs(NULL, CL_DEVICE_TYPE_GPU, 1, &devices[1], NULL);

// Create a context include two devices
cl_context ctx;
ctx = clCreateContext(0, 2, devices, NULL, NULL);

// Create queues to each device
cl_command_queue queue_cpu, queue_gpu;
queue_cpu = clCreateCommandQueueWithProperties(context, devices[0], 0, NULL);
queue_gpu = clCreateCommandQueueWithProperties(context, devices[1], 0, NULL);
\end{codebox}

代码清单5.2 使用同一个上下文为两个不同设备创建命令队列

OpenCL下的多设备编程，可以总结为以下的两种情况(使用异构设备进行并行编程)：

\begin{itemize}
\item 流水执行：两个或多个设备以流水方式工作，这样就需要设备间互相等待结果，如图5.2所示
\item 单独执行：这种方式就是每个设备各自做各自的任务，每个设备间的任务并无相关性，如图5.3所示
\end{itemize}

代码清单5.3中，等待列表执行的顺序为：CPU需要等待GPU上的内核全部执行完成，才能执行自己的内核(图5.2所示)。

\begin{codebox}[cpp]
cl_event event_cpu, event_gpu;

// Starts as soon as enqueued
err = clEnqueueNDRangeKernel(queue_gpu, kernel_gpu, 2, NULL, global, local, 0, NULL, &event_gpu);

// Starts after event_gpu is on CL_COMPLETE
err = clEnqueueNDRangeKernel(queue_cpu, kernel_cpu, 2, NULL, global, local, 0, NULL, &event_cpu);
\end{codebox}

代码清单5.3 使用流水方式进行多设备合作。CPU端的入队命令要等到GPU上的内核完成后才能执行。

\myGraphic{0.85}{images/chapter5/5-2.png}{图5.2 多个设备以流水方式处理同一份数据。CPU队列会等待GPU内核执行完毕。}

\myGraphic{0.85}{images/chapter5/5-3.png}{图5.3 多个设备并行工作。此时两个GPU不使用相同的缓冲区，各自独立执行；CPU队列要等两个GPU设备都完成后才继续。}

代码清单5.4展示了并行执行在不同设备上的内核。图5.4展示了两个GPU设备独立执行自己的内核。CPU端的入队命令要等到GPU上的内核完成后才能执行。

\begin{codebox}[cpp]
cl_event events_gpu[2];

// Both of the GPU devices can execute concurrently as soon as they have
// their respective data since they have no events in their wait-lists
err = clEnqueueNDRangeKernel(queue_gpu_0, kernel_gpu, 2, NULL, global, local, 0, NULL, &event_gpu[0]);
err = clEnqueueNDRangeKernel(queue_gpu_1, kerenl_gpu, 2, NULL, global, local, 0, NULL, &event_gpu[1]);

// The CPU will wait until both GPUs have finished executing their kernels.
// This requires two events in the CPU's wait-list
err = clEnqueueNDRangeKernel(queue_cpu, kernel_cpu, 2, NULL, global, local, 2, events_gpu, NULL);
\end{codebox}

代码清单5.4 多设备并行。CPU端的入队命令要等到GPU上的内核完成后才能执行。



% ---- 5.3-chinese.md ----
\section{内核执行域:工作项、工作组和NDRange}

OpenCL的执行以内核为中心。内核代码用来表示一个内核线程所要做的所有事情，其代码由OpenCL C语言完成。内核代码第一眼看去和C函数很类似：OpenCL C语言中，内核可以看成一个C函数。OpenCL内核具有一组参数列表，还有局部变量(类似Pthread中线程的局部变量)，以及标准的控制流。OpenCL内核的并行性，使其区别于C函数。我们在第3章仅用OpenCL工作项，并行化处理了一维空间中的大量数据。本节，我们将继续延伸讨论OpenCL执行模型中，工作项、工作组，以及NDRange的层级结构。

内核的执行需要运行时提供对应的启动/调度接口，该函数就是\texttt{clEnqueueNDRangeKernel()}。内核在调度中会产生大量工作项，共同执行内核“函数体”。使用启动接口之后，就会产生一个NDRange，其包括了内核执行的维度，以及每个维度上的工作项的数量。一个NDRange可以定义为1，2和3维，用于规划处工作项的“格子图”，工作项简单和直接的结构非常适合并行执行。将OpenCL模型映射到硬件端，每个工作项都运行在一个硬件单元上，这个单元称为“处理单元”(Processing element, PE)。OpenCL内核执行时，可能有多个工作项依次工作在同一个PE上。

内核内部调度中，每一个工作项都是独立的。OpenCL中会有意限制工作项间的同步。松散的执行模型就允许OpenCL程序去扩展设备，可用于具有超级多核心的规模化设备。可扩展的设备通常都会抽象成一种层级结构——特别是内存系统——OpenCL也为其设备的执行空间提供了一种层级结构。

为了更加灵活的支持具有大量核心的设备，OpenCL将各维度上执行的工作项等分成多个工作组。每个工作组内，工作项可以进行某种程度上的交互。OpenCL标准规定一个完整的工作组，可以并发执行在同一计算单元上。这种执行方式对于同步很重要。并发执行中的工作组内允许局部同步，不过也会对交互有所限制，以提高可扩展性。当应用中设计到需要在全局执行空间内，任务需要互相交互，那么这种OpenCL并行应用无非是低效的。使用工作组会有更高率的交互，因为一个计算单元通常映射到一个核芯上，因此工作组内交互通常都是在一块共享缓存或暂存存储器上。

通过定义更多的工作组，OpenCL内核将扩展的越来越大，并且有越来越多的线程同时在设备上执行(有更多的工作项和工作组都会在同时执行在同一设备上)。OpenCL工作项可以看做win32或POSIX线程。OpenCL层级执行模式只需要一步，因为工作项都位于工作组内，这样只需要将工作组(数量要少于工作项)映射到硬件线程的上下文中即可。这样做和单指令多数据的执行很相似，就如向量执行N操作，只使用了N个时钟周期。不过，OpenCL中子向量(工作项)可以拥有自己的程序计数器，直到同步点。用GPU来举个例子，64个工作项(英伟达为32个)可以同时被一个硬件线程处理(在SIMD单元上)，这种方式在AMD的架构中就是大家熟知的波阵面(wavefront)；在英伟达架构中，则称为线程束(warp)。即使每次执行的工作项数量被锁定，不过不同的工作项可以执行内核内不同的指令序列。这种情况会发生在内核中有分支语句时，因为If-else语句具有不同分支，所以不同的工作项会对条件状态进行评估后执行。不过工作项也有可能在同一波或束中，走进不同的分支中，这时硬件有责任将不该执行的分支结果舍弃。这种情况就是众所周知的\emph{分歧}(divergence)，这会极大影响内核执行的效率，因为工作项执行了冗余的操作，并且这些操作的结果最后都要舍弃。

这样的执行模型，就需要所有的工作项具有自己的程序计数器，这种方式要比显式使用SIMD多媒体扩展流指令在x86的处理上简单许多。因为SIMD的执行方式，需要注意给定的设备，了解对应设备支持的SSE版本，OpenCL的工作组尺寸通常是SIMD最大向量宽度的倍数。工作组尺寸查询API为\texttt{clGetKernelWorkGroupInfo()}，需要CL\_\allowbreak{}KERNEL\_\allowbreak{}PREFERRED\_\allowbreak{}WORK\_\allowbreak{}GROUP\_\allowbreak{}SIZE\_\allowbreak{}MULTIPLE作为实参传入\textbf{param\_\allowbreak{}name}中。

OpenCL定义了一些内置函数，可以在内核内部获取工作项在执行区域内的具体位置。这些函数通常都有一个维度参数，这里定义为uint dimension。其值可以为0、1和2，可以通过维度的设置获取NDRange参数在入队时的设置：

\begin{itemize}
\item uint get\_\allowbreak{}work\_\allowbreak{}dime():返回入队时所设定的工作维度数量
\item size\_\allowbreak{}t get\_\allowbreak{}global\_\allowbreak{}size(uint dimension):返回对应维度上的全局工作项数量
\item size\_\allowbreak{}t get\_\allowbreak{}global\_\allowbreak{}id(uint dimension):返回对应维度上对应全局工作项索引
\item size\_\allowbreak{}t get\_\allowbreak{}local\_\allowbreak{}size(uint dimension):返回该维度上工作组的大小。如果内核所分配的工作项在工作组内数量不均匀，那么剩余工作组(之后章节中讨论)将返回与均匀工作项不同的值
\item size\_\allowbreak{}t get\_\allowbreak{}enqueued\_\allowbreak{}local\_\allowbreak{}size(uint dimension):返回对应维度上均匀区域内工作项数量
\item size\_\allowbreak{}t get\_\allowbreak{}num\_\allowbreak{}groups(uint dimension):返回该维度上的工作组数量，即全局大小除以入队局部大小
\item size\_\allowbreak{}t get\_\allowbreak{}group\_\allowbreak{}id(uint dimension):返回当前工作组的索引。也可通过获取该工作组的第一个元素的全局索引，然后除以工作组的大小得到当前工作组索引
\end{itemize}

这里展示一个简单的OpenCL内核实例(代码清单5.5)，其输入输出均为二维数组。图5.4展示了这段代码的执行过程。对应每个像素的工作项，调用\texttt{get\_global\_id()}都会返回不同的数值。在这个简单例子中，我们使用线性位置直接对二维数据结构进行映射。实际使用时，可能会有更加复杂的映射关系，根据输入和输出的结构，再加上算法实现，可以对数据进行处理。

\begin{codebox}[cpp]
__kernel void simpleKernel(
  __global float *a,
  __global float *b){

  int address = get_global_id(0) + get_global_id(1) * get_global_size(0);
  b[address] = a[address] * 2;
}
\end{codebox}

代码清单5.5 一个简单的内核，将输入的二维数组中的数据乘以2后，放入输出数组中

\myGraphic{0.85}{images/chapter5/5-4.png}{图5.4 执行代码清单5.5中的简单内核，图中给出了NDRange里各个工作项。}

之前版本的OpenCL中，会要求NDRange上所有维度上的工作项数量是工作组的整数倍。例如，一个NDRange的大小为800x600，其工作组的大小就不能为16x16，因为600无法整除16。实际编程时，为了程序执行的效率，通常会在执行时将不足的区域补足，这样的话NDRange的区域要比实际处理的数据的空间大一些。当我们要使用16x16的工作组大小，我们需要将NDRange的设置扩大到800x608。由于有些工作项并未映射到数据集上，这就需要对这些工作项进行处理(比如：检查工作项的索引，当发现该工作项没有对应到数据集上时，立即返回)。不过，当工作组在执行的时候进行栅栏同步时，就需要组内所有工作项都到达栅栏处，这样这种扩充的方式就无法正常工作了。为了缓解这一问题，OpenCL 2.0删除了“各维度上的工作项数量需要是工作组的倍数”这一限制。另外，OpenCL 2.0为NDRange定义了“剩余工作组”(\emph{remainder work-groups})，并且最后剩余的工作组与开发者定义的工作组大小是完全不同的。比如，800x600的NDRange和16x16的工作组大小，那剩余的工作组大小为16x8。\texttt{get\_local\_size()}返回工作组实际的大小，\texttt{get\_enqueue\_local\_size()}返回均匀工作组的大小。

OpenCL 2.0也提供了计算线性索引的内置函数，开发者不需要像使用之前版本那样自己去计算线性索引值了。线性索引对于每个工作项来说是定义明确，且唯一的标识(这个线性索引与NDRange设置的维度数量或工作组大小无关)。\texttt{get\_global\_linear\_id()}返回全局上的索引值，\texttt{get\_local\_linear\_id()}返回局部上的索引值：

\begin{itemize}
\item size\_\allowbreak{}t get\_\allowbreak{}global\_\allowbreak{}linear\_\allowbreak{}id():返回工作项的全局线性索引
\item size\_\allowbreak{}t get\_\allowbreak{}local\_\allowbreak{}linear\_\allowbreak{}id():返回工作项的组内线性索引
\end{itemize}

\subsection{同步}

使用多个执行单元的时候，使用同步机制可以约束执行的顺序。通常，OpenCL会有意的限制在不同执行单元间进行同步。这种方式会影响OpenCL的扩展性，不过OpenCL的目标是支持各种各样的异构设备。例如，执行OpenCL程序的设备会自行管理线程，比如GPU；另外也有使用操作系统管理线程的设备，比如主流的x86 CPU。不使用同步，会让程序的性能提升，不过会影响程序的正确性。对于一个x86线程来说，其会尝试访问信号量，如果该信号量不可用，则该线程阻塞；操作系统会将该线程从执行中移除，并可在几乎不受资源限制的情况下调度其他任务。对于GPU线程资源有限，GPU也会对有问题的线程使用同样的策略。例如，执行阶段移除一组线程束，而不释放其使用的资源，这就会让下一个线程束进行等待，这就如同x86线程中的信号量一样，之前的线程不释放该信号量，下一个线程则无法继续进行，从而造成程序死锁。

为了避免不必要的事情发生，OpenCL定义了阻塞式同步(例如：栅栏)，且只限于工作组内的工作项。第7章中，我们将看到OpenCL 2.0也可提供"上锁-释放"内存，用来约束使用原子操作和栅栏操作的顺序。不过，其最终目的是为了能在更加广泛的算法中，使用内存的顺序约束来保证运行顺序，而不是使用同步操作来约束执行顺序。下一节中，将讨论一些在工作组中的同步方式：栅栏和命令式同步。

\subsection{工作组栅栏}

将栅栏设置到工作组内部时，需要所有工作项都到达该栅栏才能继续下面的工作，否则任何线程都无法越过该栅栏，这是一种程序计数器级别的限制。不过，对于处于不同分支的工作项来说，其行为是未定义的：很多设备上，这将导致程序死锁，因为工作组内的有线程并未到达栅栏处，所以其他线程只能进行等待。

图5.5中展示了一个简单OpenCL同步的例子。图中，我们能看到每个工作组内都由8个工作项。OpenCL标准宽松的解释下(并未限定硬件的具体实现)，工作项在工作组内的执行并非同时。当调用\texttt{work\_group\_barrier()}函数时，会让提前到达的线程等待未到达对应位置的线程，直到所有线程到达该位置，则继续进行下面的操作。不同的工作组中完成等待的时间都是相互独立的。

内置函数\texttt{work\_group\_barrier()}的参数列表如下所示：

\begin{itemize}
\item void work\_\allowbreak{}group\_\allowbreak{}barrier(cl\_\allowbreak{}mem\_\allowbreak{}fence\_\allowbreak{}flag flags)
\item void work\_\allowbreak{}group\_\allowbreak{}barrier(cl\_\allowbreak{}mem\_\allowbreak{}fence\_\allowbreak{}flag flags, memory\_\allowbreak{}scope scope)
\end{itemize}

\myGraphic{0.85}{images/chapter5/5-5.png}{图5.5 在一次内核派发内部，执行顺序上的同步只靠栅栏在工作组内进行；全局同步由内核执行完成来保证，即完成事件意味着所有工作都已完成、内存内容符合预期。}

第7章中，我们将讨论关于scope参数的更多细节(也会对flags参数做更多的讨论)。不过，当栅栏操作完成时，flags参数就已足以决定对应的内存操作是否对组内其他工作项可见。

flags可以设置成的实参：

\begin{itemize}
\item CLK\_\allowbreak{}LOCAL\_\allowbreak{}MEM\_\allowbreak{}FENCE：需要对于局部内存的所有访问组内可见
\item CLK\_\allowbreak{}GLOBAL\_\allowbreak{}MEM\_\allowbreak{}FENCE：需要对全局内存的所有访问组内可见
\item CLK\_\allowbreak{}IMAGE\_\allowbreak{}MEM\_\allowbreak{}FENCE：需要对图像的所有访问组内可见
\end{itemize}

接下来的例程中，我们可以看到工作项将数据搬运到局部内存中，每个工作项从数组中搬运一个数，其需要工作项的操作对其他所有工作项可见。为了达到这个目的，我们会调用\texttt{work\_group\_barrier()}，并传入CLK\_\allowbreak{}LOCAL\_\allowbreak{}MEM\_\allowbreak{}FENCE作为flags的实参。

内核间的调度，所有工作都可以保证操作完全完成，以及内存保证一致。下一个内核，以同样的语义加载。假设我们入队两个内核0和1(如图5.5)，这两个内核为同一个内核对象，内核代码和API调用在代码清单5.6中展示，代码行为和图5.5一致。

这个例子中，我们可以看到，工作项只是简单的将局部内存中相邻的两个数，进行简单的操作。不过，在操作两个相邻数之前，需要保证这些用来运算的数据时可读的。

\begin{codebox}[cpp]
//------------------------------
//  Relevant host program
//------------------------------

cl_mem input = clCreateBuffer(context, CL_MEM_READ_ONLY, 32 * sizeof(float), 0, 0);

cl_mem intermediate = clCreateBuffer(context, CL_MEM_READ_WRITE, 32 * sizeof(float), 0, 0);

cl_mem output = clCreateBuffer(context, CL_MEM_WRITE_ONLY, 32 * sizeof(float), 0, 0);

clEnqueueWriteBuffer(queue, input, CL_TRUE, 0, 32 * sizeof(float), (void *)hostInput, 0, NULL, NULL);

clSetKernelArg(kernel, 0, sizeof(cl_mem), (void *)&input);
clSetKernelArg(kernel, 1, sizeof(cl_mem), (void *)&intermediate);
clSetKernelArg(kernel, 2, 8 * sizeof(float), NULL);

size_t localws[1] = {8};
size_t globalws[1] = {32};

clEnqueueNDRangeKernel(queue, kernel, 1, NULL, globalws, localws, 0, NULL, NULL);

clSetKernelArg(kernel, 0, sizeof(cl_mem), (void *)&input);
clSetKernelArg(kernel, 1, sizeof(cl_mem), (void *)&intermediate);
clSetKernelArg(kernel, 2, 8 * sizeof(float), NULL);

clEnqueueNDRangeKernel(queue, kernel, 1, NULL, globalws, localws, 0, NULL, NULL);

clEnqueueReadBuffer(queue, output, CL_TRUE, 0, 32 * sizeof(float), (void *)&hostOutput, 0, NULL, NULL);

//------------------------------
//  Kernel
//------------------------------

__kernel void simpleKernel(
  __global float *a,
  __global float *b,
  __local float *localbuf){

  // Cache data to local memory
  localbuf[get_local_id(0)] = a[get_global_id(0)];

  // Wait until all work-items have read the data and
  // it becomes visible
  work_group_barrier(CLK_LOCAL_MEM_FENCE);

  // Perform the operation and output the data
  unsigned int otherAddress = (get_local_id(0) + 1) % get_local_size(0);
  b[get_global_id(0)] = localbuf[get_local_id(0)] + localbuf[otherAddress];
}
\end{codebox}

代码清单5.6 同一个命令队列中入队两个内核

\subsection{内置工作组函数}

OpenCL C编程语言在实现的时候，有很多内置函数是基于工作组的。例如栅栏操作，该内置函数需要工作组内所有的工作项，都到达指定位置后才能继续运行。因此，当工作组中有条件块时，就需要同一组中的所有工作项都执行相同的分支。

工作组评估函数也支持各种OpenCL C内置的数据类型，例如：half, int, uint, long , ulong, flaot和double。我们可以看到，在OpenCL标准中这些属于都被gentype所替代，gentype用来表示在OpenCL C中所使用的通用数据。我们依据不同函数的功能将这些函数分成三类评估函数：

\begin{enumerate}
\item 谓词评估函数
\item 广播函数
\item 并行原语函数
\end{enumerate}

\subsection{谓词评估函数}

谓词评估函数评估工作组中的所有工作项，如果满足相关的条件则返回一个非零值。函数声明如下：

\begin{itemize}
\item int work\_\allowbreak{}group\_\allowbreak{}any(int predicate)
\item int work\_\allowbreak{}group\_\allowbreak{}all(int predicate)
\end{itemize}

评估的工作组中有一个工作项满足条件，则函数work\_\allowbreak{}group\_\allowbreak{}any()返回一个非零值。当评估的工作组中所有工作项满足条件，则函数work\_\allowbreak{}group\_\allowbreak{}all()返回一个非零值。代码清单5.7中，展示了一个work\_\allowbreak{}group\_\allowbreak{}all()函数的使用示例：

\begin{codebox}[cpp]
__kernel void compare_elements(int *input, int *output){
  int tid = get_global_id(0);
  int result = work_group_all((input[tid] > input[tid +1 ]));
  output[tid] = result;
}
\end{codebox}

代码清单5.7 谓词评估函数在OpenCL内核中的使用。

\subsection{广播函数}

广播函数是将一个工作项的数据传输给工作组内其他的工作项。函数参数在维度上进行了重载：

\begin{itemize}
\item 一维：gentype work\_\allowbreak{}group\_\allowbreak{}broadcast(gentype x, size\_\allowbreak{}t local\_\allowbreak{}id)
\item 二维：gentype work\_\allowbreak{}group\_\allowbreak{}broadcast(gentype x, size\_\allowbreak{}t local\_\allowbreak{}id\_\allowbreak{}x, size\_\allowbreak{}t local\_\allowbreak{}id\_\allowbreak{}y)
\item 三维：gentype work\_\allowbreak{}group\_\allowbreak{}broadcast(gentype x, size\_\allowbreak{}t local\_\allowbreak{}id\_\allowbreak{}x, size\_\allowbreak{}t local\_\allowbreak{}id\_\allowbreak{}y, size\_\allowbreak{}t local\_\allowbreak{}id\_\allowbreak{}z)
\end{itemize}

可以看出，函数用坐标来标定对应的工作项，然后共享该工作项的数据x。该函数将广播返回值到每个工作项上。

\subsection{并行原语函数}

OpenCL支持两类内置并行原语函数：归约和扫描。这个函数会用在很多并行应用上，这两个函数的实现都由设备供应商使用高性能代码提供，这样对于开发者来说会省去自己去优化的风险和工作量。函数声明如下：

\begin{itemize}
\item \texttt{gentype work\_group\_reduce\_\textless{}op\textgreater{}(gentype x)}
\item \texttt{gentype work\_group\_scan\_inclusive\_\textless{}op\textgreater{}(gentype x)}
\item \texttt{gentype work\_group\_scan\_exclusive\_\textless{}op\textgreater{}(gentype x)}
\end{itemize}

函数中的\texttt{\textless{}op\textgreater{}}可以替换为add，min或max。这样就可以使用该函数找到局部数组的最大值，就如下面的代码所示：

\begin{codebox}[cpp]
// float max:
max = work_group_reduce_max(local_data[get_local_id(0)]);
\end{codebox}

很多并行编程者对于前缀求和非常熟悉，前缀求和可以使用work\_\allowbreak{}group\_\allowbreak{}scan\_\allowbreak{}inclusive\_\allowbreak{}add()或work\_\allowbreak{}group\_\allowbreak{}scan\_\allowbreak{}exclusive\_\allowbreak{}add()实现。闭扫描和开扫描的不同在于，当前位置上的元素是否参与累加。闭扫描版本生成的数组会包含当前位置的元素。开扫描版本则不会包含当前位置的元素。每个工作项都可用该函数，且会在工作组内返回相关线性索引的值。

并行原语函数无法保证浮点数操作的顺序，这可能带来问题，因为浮点运算不满足结合律。



% ---- 5.4-chinese.md ----
\section{原生和内置内核}

OpenCL定义了两种不需要cl\_\allowbreak{}kernel对象的入队执行机制，就是原生内核和内置内核。原生内核和内置内核是相互独立的两个概念。原生内核提供一种机制，将标准C函数(异构设备上)入队执行。内置内核需要指定的设备，并且提供对应机制允许应用开发者启动"特殊"的硬件加速模块(有可能就是指定的设备)。

\subsection{原生内核}

原生内核是一种回调的机制，其能更简洁的集成进OpenCL的执行模型中。原生内核允许使用传统编译器去编译C标准函数(与OpenCL不同)，并将编译好的C函数在放入OpenCL的任务执行图中，由事件来触发下一个事件。原生内核可以在一个设备上入队执行，并且与OpenCL内核共享内存对象。

原生内核与OpenCL内核的区别在于设置参数方面。原生内核使用对应的API进行入队(clEnqueueNativeKernel())，将标准C函数通过指针的方式进行传入。参数列表需要连同其大小，打包传入设备。

\begin{codebox}[cpp]
cl_int
clEnqueueNativeKernel(
  cl_command_queue command_queue,
  void (CL_CALLBACK *user_func)(void *),
  void *args,
  size_t cb_args,
  cl_uint num_mem_objects,
  const cl_mem *mem_list,
  const void **args_mem_loc,
  cl_uint num_events_in_wait_list,
  const cl_event *event_wait_list,
  cl_event *event)
\end{codebox}

常规OpenCL内核可以将数组和图像作为参数传入，原生内核同样可以使用对应的数组和图像作为输入。OpenCL中向原生内核传递参数完成后，会通过一种方式进行解包。可以通过设置mem\_\allowbreak{}list的实参，向原生内核传入一个内存对象链表；args\_\allowbreak{}mem\_\allowbreak{}loc参数存储一个指针链表，用于存储解包之后的内存对象。

为了更加形象的表达，代码清单5.8中，我们使用了一个原生函数foo()，其参数链表中包含5个值，其中第0个参数存储是5，第2个存储是8；1,3参数是两个数组对象；4参数是一个图像对象(如图5.6所示)。

\begin{codebox}[cpp]
// Native function that will be enqueued to device
void foo(void *args){

  ...
}

cl_command_queueu queue = clCreateCommandQueue(...);
cl_mem buffer1 = clCreateBuffer(...);
cl_mem buffer2 = clCreateBuffer(...);
cl_mem image = clCreateImage2D(...);

void *args[5] = {(void *)5, NULL, (void *)8, NULL, NULL, NULL};

num_mem_objects = 3;
cl_mem mem_list[3] = {buffer1, buffer2, image};
void *args_mem_loc[3] = {&args[1], &args[3], &args[4]};

clEnqueueNativeKernel(queue, foo, args, sizeof(args), num_mem_objects, mem_list, args_mem_loc, 0, NULL, NULL);
\end{codebox}

代码清单5.8 将原生foo()入队

\myGraphic{0.85}{images/chapter5/5-6.png}{图5.6 clEnqueueNativeKernel()示例中内存对象与参数的对应关系。}

\subsection{内置内核}

内置内核与设备是捆绑关系，其也不会在运行时由源码进行构建。通常的内置内核会展示硬件对固定函数的加速能力，这个硬件可能是一种支持OpenCL的特殊嵌入式设备，或是自定义设备。内置内核属于OpenCL定义之外的内容，因此内置内核实现的最总解释权还在硬件供应商那里。

作为OpenCL的扩展，Intel实现了运动估计的内置内核。该扩展利用了OpenCL扩展架构，对指定领域的功能进行加速，Intel所有支持OpenCL的设备都支持该功能(进行运动估计)。



% ---- 5.5-chinese.md ----
\section{设备端排队}

OpenCL 2.x版本之前的标准，只允许命令从主机端入队。OpenCL 2.0去除了这个限制，并定义了设备端的命令队列，其允许父内核直接将子内核入队。

设备端任务队列的好处是启用了\emph{嵌套并行机制}(nested parallelism)——一个并行程序中的某个线程，再去开启多个线程[1]。嵌套并行机制常用于不确定内部某个算法要使用多少个线程的应用中。与单独使用\emph{分叉-连接}(fork-join)机制相比，单层分叉-连接机制中，线程完成任务后即退出。这两种机制的区别如图5.7所示：

\myGraphic{0.85}{images/chapter5/5-7.png}{图5.7 单层“分叉-汇合”执行模型与嵌套并行线程执行的对比。}

单独调用的方式，执行完任务后，线程的生命周期就结束了。嵌套并行将会在任务执行中间产生更多的线程，并在某个子任务完成后结束线程。

嵌套循环适用于不规则或数据驱动的循环结构。数据驱动型算法有很多，这里用一个比较常见的算法作为例子：广度优先搜索算法(BFS)。广度优先搜索算法从图中的根节点开始访问其相邻的节点。之后，访问到的节点，继续访问其附近的节点，直到所有节点都访问到。当BFS算法并行时，每个顶点会发现多少新顶点，在应用启动时并不知道。设备端入队允许开发者使用OpenCL内核实现嵌套并行，这样的方式要比在主机端入队合适的多。

总结一下设备端入队的好处：

\begin{itemize}
\item 内核可以在设备端直接入队。这样就不需要同步，或与主机的交互，并且会隐式的减少数据传输
\item 更加自然的表达算法。当算法包括递归，不规则循环结构，或其他单层并行是固定不均匀的，现在都可以在OpenCL中完美的实现
\item 更细粒度的并行调度，以及动态负载平衡。设备能更好的相应数据驱动决策，以及适应动态负载
\end{itemize}

为了让子内核入队，内核对象需要调用一个OpenCL C内置函数\texttt{enqueue\_kernel()}。这里需要注意的是，每个调用该内置函数的工作项都会入队一个子内核。该内置函数的声明如下：

\begin{codebox}[cpp]
int
enqueue_kernel(
  queue_t queue,
  kernel_enqueue_flags_t flags,
  const ndrange_t ndrange,
  void (^block)(void))
\end{codebox}

和主机端的API一样，其也需要传递一个命令队列。flags参数用来执行子内核何时开始执行。该参数有三个语义可以选择：

\begin{itemize}
\item CLK\_\allowbreak{}ENQUEUE\_\allowbreak{}FLAGS\_\allowbreak{}NO\_\allowbreak{}WAIT:子内核立刻执行
\item CLK\_\allowbreak{}ENQUEUE\_\allowbreak{}FLAGS\_\allowbreak{}WAIT\_\allowbreak{}KERNEL:子内核需要等到父内核到到ENDED点时执行。这就意味着子内核在设备上运行时，父内核已经执行完成
\item CLK\_\allowbreak{}ENQUEUE\_\allowbreak{}FLAGS\_\allowbreak{}WAIT\_\allowbreak{}WORK\_\allowbreak{}GROUP:子内核必须要等到入队的工作组执行完成后，才能执行。
\end{itemize}

需要注意的是，父内核可能不会等到子内核执行结束。一个父内核的执行状态为"完成"时，意味着其本身和其子内核都完成。如果父子内核正确无误的执行完成，那么父内核会置为CL\_\allowbreak{}COMPLETE。如果有子内核的程序计数器错误，或非正常终止，父内核的状态会置为一个错误值(一个给定的负值)。

与clEnqueueNDRangeKernel()类似，\texttt{enqueue\_kernel()}也需要使用NDRange来指定维度信息(传递给ndrange参数)。与主机端调用一样，全局偏移和工作组大小是可选的。这时需要在内核上创建ndrange\_\allowbreak{}t类型的对象来执行执行单元的配置，这里使用到了一组内置函数：

\begin{codebox}[cpp]
ndrange_t ndrange_<N>D(const size_t global_work_size[<N>])

ndrange_t ndrange_<N>D(const size_t global_work_size[<N>], const size_t global_work_size[<N>])

ndrange_t ndrange_<N>D(const size_t global_work_size[<N>], const size_t global_work_size[<N>], const size_t local_work_size_[<N>])
\end{codebox}

其中\texttt{\textless{}N\textgreater{}}可为1,2和3。例如，创建一个二维的800x600的NDRange可以使用如下方式：

\begin{codebox}[cpp]
size_t globalSize[2] = {800, 600};
ndrange_t myNdrange = ndrange_2D(globalSize);
\end{codebox}

最终，enqueue\_\allowbreak{}kernel()的最后一个参数block，其为指定入队的内核。这里指定内核的方式称为“Clang块”。下面两节中将会更加详细的对“如何利用设备入队”进行讨论，以及如何使用块语法指定一个嵌套内核。

如主机端API一样，enqueue\_\allowbreak{}kernel()会返回一个整数，代表其执行是否成功。返回CLK\_\allowbreak{}SUCCESS为成功，返回CLK\_\allowbreak{}ENQUEUE\_\allowbreak{}FAILURE则为失败。编程者想要了解失败入队的更多原因的话，需要在clBuildProgram()传入"-g"参数，或是clCompileProgram()调用会启用细粒度错误报告，会有更加具体的错误码返回，例如：CLK\_\allowbreak{}INVALID\_\allowbreak{}NDRANGE或CLK\_\allowbreak{}DEVICE\_\allowbreak{}QUEUE\_\allowbreak{}FULL。

\subsection{创建一个设备端队列}

设备端队列也需要使用clCreateCommandQueueWithProperties()在主机端进行创建。为了表明是为了设备端创建的命令队列，properties中需要传入CL\_\allowbreak{}QUEUE\_\allowbreak{}ON\_\allowbreak{}DEVICE。另外，当一个设备端队列创建之后，标准要求CL\_\allowbreak{}QUEUE\_\allowbreak{}OUT\_\allowbreak{}OF\_\allowbreak{}ORDER\_\allowbreak{}EXEC\_\allowbreak{}MODE\_\allowbreak{}ENABEL也要传入(使用乱序方式)，OpenCL 2.0中设备端队列被定义为乱序队列。这时命令队列可以通过内核的参数传入内核内，其对应的就是内核中的queue\_\allowbreak{}t类型。代码清单5.9中，展示了一个带有队列参数的内核。

\begin{codebox}[cpp]
// -----------------------------
// Relevant host program
// -----------------------------

// Specify the queue properties
cl_command_queue_properties properties =
  CL_QUEUE_ON_DEVICE |
  CL_QUEUE_OUT_OF_ORDER_EXEC_MODE_ENABLE;

// Create the device-side command-queue
cl_command_queue device_queue;
device_queue = clCreateCommandQueueWithProperties(
  context,
  device,
  &proerties,
  NULL);

...

clSetKernelArg(kernel, 0, sizeof(cl_command)，&device_queue);

...

// -----------------------------
// Kernel
// -----------------------------

__kernel
void foo(queue_t myQueue, ...){
  ...
}
\end{codebox}

代码清单5.9 将设备端队列作为参数设置到内核中

有另外一个可选的配置，CL\_\allowbreak{}QUEUE\_\allowbreak{}ON\_\allowbreak{}DEVICE\_\allowbreak{}DEFAULT可传入clCreateCommandQueueWithProperties中，这样产生的队列默认为设备端可用队列。对于编程者来说，这样简化了许多工作，因为默认队列就可以在相关的内核使用内置函数，所以就不需要再将命令队列作为参数传入内核中了(可以通过内置函数get\_\allowbreak{}default\_\allowbreak{}queue()获取队列对象)。

\subsection{入队设备端内核}

当使用主机端API(clEnqueueNDRangeKernel())入队内核执行命令之前，需要对内核进行参数的设置。设备端使用enqueue\_\allowbreak{}kernel()将内核命令入队前，是没有设置参数的过程。不过，内核的参数还是要设置，那该如何是好呢？为了正确的执行内核，OpenCL选择使用Clang块语法进行参数设置。

Clang块在OpenCL标准中，作为一种方式对内核进行封装，并且Clang块可以进行参数设置，能让子内核顺利入队。Clang块是一种传递代码和作用域的方法。其语法与闭包函数和匿名函数类似。块类型使用到一个结果类型和一系列参数类型(类似lambda表达式)。这种语法让“块”看起来更像是一个函数类型声明。“\textasciicircum{}”操作用来声明一个块变量(block variable)(该块用来引用内核)，或是用来表明作用域的开始(使用内核代码直接进行声明)。

代码清单5.10中的简单例子，用来表明块语法的使用方式。

\begin{codeboxt}[cpp]{代码清单5.10 一个用来展示Clang块语法是如何使用的例子}
__kerenl
void child0_kernel(){
  printf("Child0: Hello, world!\n");
}

void child1_kernel(){
  printf("Child1: Hello, world!\n");
}

__kernel
void parent_kernel(){
  kernel_enqueue_flags_t child_flags = CLK_ENQUEUE_FLAGS_NO_WAIT;
  ndrange_t child_ndrange = ndrange_1D(1);

  // Enqueue the child kernel by creating a block variable
  enqueue_kernel(get_default_queue(), child_flags, child_ndrange, ^{child0_kernel();});

  // Create block variable
  void (^child1_kernel_block)(void) = ^{child1_kernel()};

  // Enqueue kernel from block variable
  enqueue_kernel(get_default_queue(), child_flags. child_ndrange, child1_kernel_block);

  // Enqueue kernel from a block literal
  // The block literal is bound by "^{" and "}"
  enqueue_kernel(get_default_queue(), child_flags, child_ndrange, ^{printf("Child2: Hello, world!\n");});
}
\end{codeboxt}

代码清单5.10中展示了三种设备端入队的方式。不过，这里的内核都没有参数。块语法是支持传递参数的。这里将第3章的向量相加的内核借用过来：

\begin{codebox}[cpp]
__kernel
void vecadd(
  __global int *A,
  __global int *B,
  __global int *C){

  int idx = get_global_id(0);

  C[idx] = A[idx] + B[idx];
}
\end{codebox}

可以用这个内核来进行参数传递，我们可以将这个内核作为一个子内核进行入队。代码清单5.11中展示了如何使用块变量进行参数传递，代码清单5.12中使用了块函数进行参数传递。注意代码清单5.11中的参数传递有些类似于标准函数的调用。不过，当我们使用块函数的方式时，不需要显式的传递参数。编译器会为作用域创建一个新的函数域。全局变量可以使用绑定的方式，私有和局部数据就必须进行拷贝了。注意指向私有或局部地址的指针是非法的，因为当程序运行到工作组或工作项之外时，这些指针就会失效。不过，OpenCL也支持为子内核开辟局部内存，我们将会在后面讨论这种机制。因为内核函数总是返回void，声明块时不需要显式的定义。

\begin{codeboxt}[cpp]{代码清单5.11 使用块语法传递参数}
__kernel
void child_vecadd(
  __global int *A,
  __global int *B,
  __global int *C){

  int idx = get_gloabl_id(0);

  C[idx] = A[idx] + B[idx];
}

__kernel
void parent_vecadd(
  __global int *A,
  __global int *B,
  __global int *C){

  kernel_enqueue_flags_t child_flags = CLK_ENQUEUE_FLAGS_NO_WAIT;
  ndrange_t child_ndrange = ndrange_1D(get_global_size(0));

  // Only enqueue one child kernel
  if (get_global_id(0) == 0){

    enqueue_kerenl(
      get_default_queueu(),
      child_flags,
      child_ndrange,
      ^{child_vecadd(A, B, C);}); // Pass arguments to child
  }
}
\end{codeboxt}

\begin{codeboxt}[cpp]{代码清单5.12 使用作用域的方式}
__kernel
void parent_vecadd(
  __global int *A,
  __global int *B,
  __global int *C){

  kernel_enqueue_flags_t child_flags = CLK_ENQUEUE_FLAGS_NO_WAIT;
  ndrange_t child_ndrange = ndrange_1D(get_global_size(0));

  // Only enqueue one child kernel
  if (get_global_id(0) == 0){

    // Enqueue kernel from block literal
    enqueue_kerenl(
      get_default_queueu(),
      child_flags,
      child_ndrange,
      ^{int idx = get_global_id(0);
        C[idx] = A[idx] + B[idx];});
  }
}
\end{codeboxt}

\textbf{动态局部内存}

使用主机端API设置内核参数，需要动态分配局部内存时，只需要使用clSetKernelArg()向内核该参数传递为NULL就好。设备端没有类似的设置参数的机制，enqueue\_\allowbreak{}kernel()有重载的函数：

\begin{codebox}[cpp]
int
enqueue_kerenl(
  queue_t queue,
  kernel_enqueue_flags_t flags,
  const ndrange_t ndrange,
  void (^block)(local void *, ...),
  uint size0, ...)
\end{codebox}

该函数可用来创建局部内存指针，标准中块可以是一个可变长参数的表达式(一个可以接受可变长参数的函数)。其中每个参数的类型必须是\texttt{local void *}。注意，声明中，函数列表可以被void类型替代。enqueue\_\allowbreak{}kernel()函数同样也是可变长的，末尾提供的数值是用来表示每个局部素组的大小。代码清单5.13对向量相加的例子进行修改，用来展示如何使用块语法，进行动态局部内存的分配。

\begin{codeboxt}[cpp]{代码清单5.13 如何动态的为子内核动态分配局部内存}
// When a kernel has been defined like this, then it can be
// enqueued from the host as well as from the device
__kernel
void child_vecadd(
  __global int *A,
  __global int *B,
  __global int *C,
  __local int *localA,
  __local int *localB,
  __local int *localC){

  int idex = get_global_id(0);
  int local_idx = get_local_id(0);

  local_A[local_idx] = A[idx];
  local_B[local_idx] = B[idx];
  local_C[local_idx] = local_A[local_idx] + local_B[local_idx];
  C[idx] = local_C[local_idx];
}

__kernel
void parent_vecadd(
  __global int *A,
  __global int *B,
  __global int *C){

  kernel_enqueue_flags_t child_flags = CLK_ENQUEUE_FLAGS_NO_WAIT;
  ndrange_t child_ndrange = ndrange_1D(get_global_size(0));

  int local_A_mem_size = sizeof(int) * 1;
  int local_B_mem_size = sizeof(int) * 1;
  int local_C_mem_size = sizeof(int) * 1;

  // Define a block with local memeory for each
  // local memory argument of the kernel
  void (^child_vecadd_blk)(
    local int *,
    local int *,
    local int *) =
    ^(local int *local_A,
      local int *local_B,
      local int *local_C){

      child_vecadd(A, B, C, local_A, local_B, local_C);
    };

  // Only enqueue one child kernel
  if (get_global_id(0) == 0){
    // Variadic enqueue_kernel function takes in local
    // memory size of each argument in block
    enqueue_kernel(
      get_default_queue(),
      child_flags,
      child_ndrange,
      child_vecadd_blk,
      local_A_mem_size,
      local_B_mem_size,
      lcoal_C_mem_size);
  }
}
\end{codeboxt}

\textbf{使用事件强制依赖}

介绍设备端命令队列时就曾说过，设备端命令队列可以乱序的执行命令。这就暗示着设备端需要提供严格的依赖顺序。在主机端将内核直接入队，事件对象可以很好的指定依赖顺序。这里再次提供一个enqueue\_\allowbreak{}kernel()的另一个重载版本。

\begin{codebox}[cpp]
int
enqueue_kenrel(
  queue_t queue,
  kernel_enqueue_flags_t flags,
  const ndrange_t ndrange,
  uint num_events_in_wait_list,
  const clk_event_t *event_wait_list,
  clk_event_t *event_ret,
  void (^block)(void))
\end{codebox}

读者需要注意的是增加的三个参数：num\_\allowbreak{}events\_\allowbreak{}in\_\allowbreak{}wait\_\allowbreak{}list，event\_\allowbreak{}wait\_\allowbreak{}list和event\_\allowbreak{}ret。与主机端API中事件相关的参数用法一样。

事件版本当然也有局部内存的重载版本：

\begin{codebox}[cpp]
int
enqueue_kenrel(
  queue_t queue,
  kernel_enqueue_flags_t flags,
  const ndrange_t ndrange,
  uint num_events_in_wait_list,
  const clk_event_t *event_wait_list,
  clk_event_t *event_ret,
  void (^block)(local void *, ...),
  uint size0, ...)
\end{codebox}

\begin{codeboxt}[cpp]{代码清单5.14 设备端入队时使用事件指定依赖关系}
__kernel
void child0_kernel(){
  printf("Child0: I will run first.\n");
}

__kernel
void child1_kernel(){
  printf("Child1: I will run second.\n");
}

__kernel
void parent_kernel(){
  kernel_enqueue_flags_t child_flags = CLK_ENQUEUE_FLAGS_NO_WAIT;
  ndrange_t child_ndrange = ndrange_1D(1);

  clk_event event;

  // Enqueue a kernel and initialize an event
  enqueue_kernel(
    get_default_queue(),
    child_flags,
    child_ndrange,
    0,
    NULL,
    &event,
    ^{child0_kernel();});

  // Pass the event as a dependency between the kernels
  enqueue_kernel(
    get_default_queue(),
    child_flags,
    child_ndrange,
    1,
    &event,
    NULL,
    ^{child1_kernel();});

  // Release the event. In this case, the event will be released
  // after the dependency is satisfied
  // (second kernel is ready to execute)
  release_event(event);
}
\end{codeboxt}

\myRefRule

\begin{myRefs}

[1] J.Reinders. Intel Threading Building Blocks: Outfiting C++ for Multi-Core Processor Parallelism. O'Reilly Media, Inc., Sebastopol. 2007.

\end{myRefs}



% ---- 5.6-chinese.md ----
\section{本章总结}

本章中我们讨论了很多有关运行时和执行模型的话题。运行时提供基于任务的执行模式，其底层是基于OpenCL的入队模型和事件对象提供的依赖机制。事件对象能用来分析OpenCL的命令，以及“用户定义的相关回调命令”的性能。内核的执行域是一组层级结构，其包含工作项，工作组，以及NDRange。OpenCL 2.0中加入的一些新功能，能让内核在设备端入队。我们也了解了如何在内核中使用块语法。同样，执行顺序和基本的同步方式也在本章中有所涉及，后续的章节中将描述主机端和设备端的内存模型，以及更加深入的了解执行单元的交互和同步。


